\documentclass[11pt]{article}
% Shared local declarations. Document roots use the repository's business-admin family.
\usepackage{business-admin}
\usepackage{amsmath,amssymb,tabularx,longtable,array,booktabs,enumitem}
\usepackage{ragged2e}
\newcolumntype{Y}{>{\RaggedRight\arraybackslash}X}
\newcommand{\SuiteMeta}[2]{\def\DocTitle{#1}\def\DocSubtitle{#2}\def\DocVersion{1.0}\def\DocStatus{Educational and operational reference}\def\DocOwner{Parliamentary Procedure and Governance Suite}\def\DocSummary{Independently authored reference and working materials, based on RONR 12th edition.}}
\newcommand{\SuiteOpen}{\makedoctitle\vspace{-0.5em}\noindent\rule{\linewidth}{0.4pt}\par\medskip}
\newcommand{\RonrAnchor}[1]{\par\smallskip{\footnotesize\textbf{Primary reading:} Robert's Rules of Order Newly Revised, 12th ed., #1. Original synthesis; not official interpretation.}\par}
\newcommand{\WriteLine}{\par\noindent\rule{\linewidth}{0.35pt}\par\vspace{0.5cm}}
\newcommand{\WriteLines}[1]{\par\noindent\rule{\linewidth}{0.35pt}\par\vspace{0.35cm}\noindent\rule{\linewidth}{0.35pt}\par\vspace{0.35cm}\noindent\rule{\linewidth}{0.35pt}\par\vspace{0.25cm}}
\newcommand{\Field}[1]{\par\medskip\noindent\textbf{#1}\par\WriteLines{3}}
\newcommand{\GovCheckbox}{\(\square\)}
\setlength{\emergencystretch}{3em}

\SuiteMeta{Laboratory: Vote Threshold and Abstention Clinic}{Scenario | simulation, evidence and assessment}
\begin{document}
\SuiteOpen
\section{Learning target}
Produce a defensible sequence of parliamentary decisions and corresponding organizational evidence. Cite the 12th-edition rule before using it to resolve a disagreement.

\section{Initial conditions}
Twenty-one members are present. On a motion requiring a majority of votes cast, eleven vote yes, nine no and one abstains; then assume the same count under a two-thirds requirement.

\section{Roles and materials}
\begin{itemize}[leftmargin=*,itemsep=4pt]
\item Presiding officer, secretary, proposer, objector or challenger as applicable; independent observer.
\item Fictional bylaws excerpt and adopted special rules; printed motion cards and tally sheet.
\item The corresponding template forms and blank copies in this suite.
\end{itemize}

\section{Procedure}
\begin{enumerate}[leftmargin=*,itemsep=4pt]
\item Compute minimum yes count by basis and write announcements.
\item Before any vote, check valid meeting, membership, quorum, motion class, second, debate and threshold.
\item At the procedural branch, explain the governing rule and offer the correct next action.
\item Announce the outcome and preserve exact decision language, tally basis and required evidence.
\item Have an independent observer reproduce the sequence from the record alone.
\end{enumerate}

\section{Expected reasoning}
Majority passes (11 of 20); two-thirds fails (minimum 14 of 20).

If organizational bylaws differ from general RONR procedures, document the conflict and follow controlling higher authority. The scenario assumes no additional statutes or special rules unless supplied.

\section{Required submissions}
\begin{itemize}[leftmargin=*,itemsep=4pt]
\item Authority and quorum validation sheet.
\item Procedural event log including the pending question after each action.
\item Completed motion or teller form with counts and exact threshold.
\item Action-only draft minutes excerpt and discrepancy report.
\end{itemize}

\section{Evaluation rubric}
Accuracy of cited rule (30\%); motion sequencing and fairness (30\%); quantitative or authority correctness (20\%); evidence completeness and clarity (20\%). No credit for invented bylaws or undocumented votes.
\end{document}
